\subsection{Node Property Flow (NPF) Package}

The Node Property Flow (NPF) Package computes intercell conductances from the
hydraulic conductivity input and cell geometry. The single most common API
modification for this package is changing intercell conductance, which is done
by setting the derived \texttt{CONDSAT} array rather than the hydraulic
conductivity input. Under the standard (non-XT3D) formulation, the hydraulic
conductivity arrays are read only during Allocate and Read, where they are used
to compute \texttt{CONDSAT}; changing them after that point has no effect.

\begin{longtable}{p{0.15\textwidth} p{0.15\textwidth} p{0.60\textwidth}}
\caption{Memory-manager variables in the NPF Package that are candidates for
modification through the API. The memory address of each variable is
\texttt{MODELNAME/NPF/VARNAME}.} \label{table:api-gwf-npf} \\

\hline
\textbf{Variable} & \textbf{Class} & \textbf{Guidance} \\
\hline
\endfirsthead

\hline
\textbf{Variable} & \textbf{Class} & \textbf{Guidance} \\
\hline
\endhead

\texttt{CONDSAT} & Derived variable &
Saturated intercell conductance, dimensioned by the number of symmetric
connections (NJAS). Computed once during Allocate and Read from the hydraulic
conductivity arrays, cell geometry, and \texttt{ICELLTYPE}, then used directly
in every formulate step. Set \texttt{CONDSAT} to change intercell conductance at
run time. \emph{Caveats:} the array is allocated with size zero when the XT3D
option is active (conductance is then built by XT3D from the conductivity
arrays). It is recomputed every time step when time-varying hydraulic
conductivity (TVK) or the Viscosity (VSC) Package is active, which overwrites any
externally set values. \\

\texttt{K11} & Inert input &
Hydraulic conductivity along the major axis, dimensioned by the number of cells
(NODES). Used only to compute \texttt{CONDSAT} during Allocate and Read; setting
it afterward has no effect under the standard formulation---set \texttt{CONDSAT}
instead. \emph{Exceptions:} when XT3D is active the conductance is built from the
conductivity arrays during formulation, so \texttt{K11} becomes influential; when
TVK is active \texttt{K11} is re-read each time step and overwrites API changes
(class becomes Overwritten). \\

\texttt{K22} & Inert input &
Hydraulic conductivity along the second axis, dimensioned by NODES (present when
the \texttt{K22} input is specified). Same guidance and exceptions as
\texttt{K11}. \\

\texttt{K33} & Inert input &
Hydraulic conductivity along the third (commonly vertical) axis, dimensioned by
NODES (present when the \texttt{K33} input is specified). Same guidance and
exceptions as \texttt{K11}. \\

\texttt{ANGLE1}, \texttt{ANGLE2}, \texttt{ANGLE3} & Inert input &
Rotation angles of the hydraulic conductivity ellipsoid, dimensioned by NODES
(present only when specified). Used to compute \texttt{CONDSAT} (or the XT3D
coefficients) during Allocate and Read; set \texttt{CONDSAT} instead under the
standard formulation. \\

\texttt{SAT} & Overwritten &
Saturated fraction of the cell (0 to 1), dimensioned by NODES. Recomputed from
the current head at every formulate step for convertible cells, so values set
through the API are overwritten. To influence saturation, change the head
(\texttt{X}) or the cell geometry instead. \\

\texttt{ICELLTYPE} & Structural / read-only &
Cell convertibility flag (0 for confined, nonzero for convertible), dimensioned
by NODES. Read during Allocate and Read to compute \texttt{CONDSAT} and consulted
during formulation. Changing it at run time is not supported and can produce
inconsistent conductance and storage; treat as read-only. \\

\hline
\end{longtable}
